% needed for custom switches
\RequirePackage{etoolbox}

% checkmode switch:
% show all overfull/underfull boxes, show unknown hyphenations
\newtoggle{checkmode}
\toggletrue{checkmode}

% debugmode switch:
% show all boxes, glues, and kerning info
\newtoggle{debugmode}
\togglefalse{debugmode}

% warns about outdated packages and missing caption declarations
\RequirePackage[
  l2tabu,
  orthodox,
]{nag}

% set main document class and options
\documentclass[
  % choose font size (PO: should be between 12pt and 14pt)
  fontsize=12pt,
  % let typearea calculate appropriate page margins
  DIV=calc,
  % don't use KOMA-Script's default "big" headings
  headings=normal,
]{scrbook}

% check mode
\iftoggle{checkmode}{
  % show overfull boxes
  \overfullrule=1mm
  % show underfull vboxes
  \directlua{dofile("detect_underfull.lua")}
  % show non-whitelisted hyphenations
  \usepackage[mark]{lua-check-hyphen}
  % whitelist is in hyphenation_whitelist.txt:
  % one word per line with dashes (-) indicating the places where
  % hyphenation is whitelisted
  \LuaCheckHyphen{whitelist=hyphenation_whitelist.txt}
}{}

% debug mode
\iftoggle{debugmode}{
  % show boxes, glues, and kerning
  \usepackage{lua-visual-debug}
}{}

% language-specific things (hyphenation, ...)
\usepackage[ngerman,american]{babel}

% biblatex package wants to have csquotes (otherwise: warning)
\usepackage{csquotes}

% dummy texts
\usepackage[math]{blindtext}

% basic AMS math commands
\usepackage{amsmath}

% micro-typographic adjustments
% (final: don't deactivate extensions due to document class draft option)
\usepackage[final]{microtype}

% Unicode umlauts
\usepackage[utf8]{luainputenc}

% need T1 font encoding for Charter,
% otherwise there will be "undefined font shape" warnings
\usepackage[T1]{fontenc}

% use Bitstream Charter as main font
\usepackage[bitstream-charter]{mathdesign}

% don't use another, sans-serif font for headings
\setkomafont{disposition}{\normalcolor\bfseries}

% graphics
\usepackage{graphicx}

% X column type for tables (filling rest of line)
\usepackage{tabularx}

% line spacing (one and a half lines)
\usepackage[onehalfspacing]{setspace}

% location of graphics files
\graphicspath{{../gfx/}}

% chapter heading: make number bigger
\renewcommand*{\chapterformat}{%
  \mbox{%
    \chapappifchapterprefix{\nobreakspace}%
    \scalebox{3.5}{\thechapter\autodot}%
    \IfUsePrefixLine{}{\enskip}%
  }%
}

% chapter heading: align number and text at the baseline of the last line
% of the text, add some more space between number and text
\renewcommand*{\chapterlinesformat}[3]{%
  \begin{tabularx}{\textwidth}{@{}l@{}X@{}}%
    % chapter number
    #2%
    % add hspace between number and text, but only if \chapter (not \chapter*)
    % has been used (i.e., if the chapter number #2 is not empty)
    \if\relax\detokenize{#2}\relax\else\hspace*{5mm}\fi&%
    % text of heading, [b] is for alignment, \linewidth is width of X cell,
    % \raggedchapter for not justifying headings
    \parbox[b]{\linewidth}{\raggedchapter#3}%
  \end{tabularx}%
}

% PDF links
\usepackage{hyperref}

% BibLaTeX
\usepackage{biblatex}

% location of *.bib file
\addbibresource{../bib/bibliography.bib}

% named references
\usepackage{cleveref}

% meta-data variables
\newcommand*{\thetitle}{%
  B-Splines for Sparse Grids:\texorpdfstring{\\}{}
  Algorithms and Application to
  Higher-Dimensional Optimization%
}
\newcommand*{\theauthor}{Julian Valentin}
\newcommand*{\thedate}{TODO(date)}

% set up hyperref
\hypersetup{
  pdftitle={\thetitle},
  pdfauthor={\theauthor},
  pdfcreator={LaTeX, KOMA-Script, hyperref},
}

% change bibliography heading (default: "Bibliography")
\DefineBibliographyStrings{english}{bibliography={References}}

\begin{document}
  % title page
  \begin{titlepage}
  \begin{spacing}{1}
    \begin{center}
      \begin{otherlanguage}{ngerman}
        % no indentation of paragraphs
        \setlength{\parindent}{0pt}
        
        {\bfseries\huge\thetitle\par}
        
        \vfill
        
        Vom Stuttgarter Zentrum für Simulationswissenschaften der\\
        Universität~Stuttgart zur Erlangung der Würde eines Doktors der\\
        Naturwissenschaften (Dr.~rer.~nat.) genehmigte Abhandlung
        
        \vfill
        
        Vorgelegt von
        
        {\bfseries\Large\theauthor\par}
        
        aus \thebirthplace
        
        \vfill
        
        \begin{tabular}{ll}
          Hauptberichter:&
          \theadvisor\\[0.5em]
          Mitberichter:&
          \theexamineri\\
          &\theexaminerii\\[1em]
          \multicolumn{2}{l}{%
            Tag der mündlichen Prüfung:\quad%
            \thedefensedate%
          }
        \end{tabular}
        
        \vfill
        
        \includegraphics[scale=1.1]{logoUniversityGerman}
        
        \vspace{2em}
        
        Vorgelegt an der \theuniversity
        
        \vspace{1em}
        
        \theinstitute{} der \theuniversity
        
        \vspace{1em}
        
        \theyear
      \end{otherlanguage}
    \end{center}
  \end{spacing}
\end{titlepage}

% don't display page number
\thispagestyle{empty}

{%
  % no indentation of paragraphs
  \setlength{\parindent}{0pt}%
  \small
  
  \begin{center}
    \includegraphics[scale=1.1]{logoUniversityEnglish}%
    
    \vspace{1em}
    
    Submitted to the University of Stuttgart
  \end{center}
  
  \begin{tabular}{@{}p{0.57\textwidth}@{}p{0.43\textwidth}@{}}
    \emph{Involved institutions and departments:}%
    \vspace{0.6mm}\newline%
    Cluster of Excellence in Simulation Technology%
    \vspace{0.6mm}\newline%
    Institute for Parallel and Distributed Systems%
    \vspace{0.6mm}\newline%
    Department for Simulation Software Engineering%
    \vspace{0.6mm}\newline%
    Department for Simulation of Large Systems&
    \raisebox{-0.5\height}{%
      \includegraphics[scale=1.0]{logoSimTech}%
    }%
    \qquad%
    \raisebox{-0.5\height}{%
      \includegraphics[scale=1.2]{logoIPVS}%
    }%
    \vspace{2mm}\newline%
    \raisebox{-0.3888\height}{%
      \includegraphics[scale=1.0]{logoSSE}%
    }%
    \qquad%
    \raisebox{-0.5\height}{%
      \includegraphics[height=19mm]{logoSGS.png}%
    }
  \end{tabular}
  
  \vfill
  
  Compiled as version \texttt{v\compileCounter{}}
  on \currentTimeLong.\\
  \lefthphantom{%
    Committed as \texttt{\gitCommitHash{}}%
  }{%
    Compiled as version \texttt{v\compileCounter{}}%
  }
  on \gitCommitTimeLong.
  
  \vspace{1em}
  
  Typeset using \LaTeX{} and cover design by the author.
  
  \vspace{1em}
  
  \begin{tabular}{%
    @{}p{0.2\textwidth}@{}p{0.8\textwidth}@{}%
  }
    \includegraphics[height=10mm]{licenseBadge}&
    \raisebox{3.8mm}{%
      \parbox{\linewidth}{%
        Copyright \copyright{} \theyear{} \theauthor{}.
        This work is licensed under the
        \term{%
          \href{https://creativecommons.org/licenses/by-sa/4.0/}{%
            Creative Commons Attribution-ShareAlike 4.0
            International License%
          }%
        }.%
      }%
    }
  \end{tabular}
  
  \vspace{1em}
  
  Although this thesis was written with uttermost care,
  it cannot be ruled out that it contains errors.
  Please send any corrections and mistakes to
  \href{mailto:thesis@bsplines.org}{\texttt{thesis@bsplines.org}}.
}

\cleardoublepage

  
  % dedication page
  \cleardoublepage
  % don't display page number
\thispagestyle{empty}

\vspace*{\fill}

\begin{center}
  \includegraphics[width=0.6\textwidth]{dedicationQuote}
  
  \begin{minipage}{0.6\textwidth}%
    \begin{flushright}
      \small--- Carl de~Boor \cite{Boor16Comment}
    \end{flushright}
  \end{minipage}
\end{center}

\vspace*{\fill}

\begin{center}
  \includegraphics{logoBSplines}%
\end{center}

\vspace*{\fill}

\cleardoublepage

  
  % table of contents
  \tableofcontents
  
  % abstract
  \addchap{Abstract/\foreignlanguage{ngerman}{Kurzzusammenfassung}}

\printornamentsfalse

\section*{Abstract}

In simulation technology, computationally expensive objective functions
are often replaced by cheap surrogates,
which can be obtained by interpolation.
Full grid interpolation methods suffer from the
so-called curse of dimensionality,
rendering them infeasible if the parameter domain of the function
is higher-dimensional (four or more parameters).
Sparse grids constitute a discretization method that does not suffer from the
curse, while the approximation quality deteriorates only insignificantly.
However, conventional basis functions such as piecewise linear functions
are not smooth (continuously differentiable).
Hence, these basis functions are unsuitable for applications
in which gradients are required.
One example for such an application is gradient-based optimization,
in which the availability of gradients greatly improves the speed of
convergence and the accuracy of results.

This thesis demonstrates that hierarchical B-splines on sparse grids are
well-suited for obtaining smooth interpolants for higher dimensionalities.
The thesis is organized in two main parts:
In the first part, we derive new B-spline bases on sparse grids and study
their implications on theory and algorithms.
In the second part, we consider three real-world applications in optimization:
topology optimization, biomechanical continuum-mechanics, and
dynamic portfolio choice models in finance.
The results reveal that the optimization problems of these applications
can be solved accurately and efficiently with hierarchical B-splines on
sparse grids.
% 209 words

\newpage

\begin{otherlanguage}{ngerman}
  \section*{Kurzzusammenfassung}
  
  In der Simulationstechnologie werden zeitaufwendige Zielfunktionen
  oft durch einfache Surrogate ersetzt, die durch Interpolation
  gewonnen werden können.
  Vollgitter-Interpolationsmethoden leiden unter dem
  sogenannten Fluch der Dimensionalität,
  der sie unbrauchbar macht, falls der Parameterbereich der Funktion
  höherdimensional ist (vier oder mehr Parameter).
  Dünne Gitter sind eine Diskretisierungsmethode, die nicht unter
  dem Fluch leidet, aber die Approximationsqualität nur leicht verschlechtert.
  Leider sind konventionelle Basisfunktionen wie die stückweise
  lineare Funktionen nicht glatt (stetig differenzierbar).
  Daher sind sie für Anwendungen ungeeignet, in denen Gradienten
  benötigt werden.
  Ein Beispiel für eine solche Anwendung ist gradientenbasierte Optimierung,
  in der die Verfügbarkeit von Gradienten die Konvergenzgeschwindigkeit und
  die Ergebnisgenauigkeit deutlich verbessert.
  
  Diese Dissertation demonstriert, dass hierarchische B-Splines auf
  dünnen Gittern gut geeignet sind,
  um glatte Interpolierende für höhere Dimensionalitäten zu erhalten.
  Die Dissertation ist in zwei Hauptbereiche gegliedert:
  Der erste Teil leitet neue B-Spline-Basen auf dünnen Gittern her und
  untersucht ihre Implikationen bezüglich Theorie und Algorithmen.
  Der zweite Teil behandelt drei Realwelt-Anwendungen aus der Optimierung:
  Topologieoptimierung, biomechanische Kontinuumsmechanik und
  Modelle der dynamischen Portfolio-Wahl in der Finanzmathematik.
  Die Ergebnisse zeigen, dass die Optimierungsprobleme dieser
  Anwendungen durch hierarchische B-Splines auf dünnen Gittern
  genau und effizient gelöst werden können.
  % 188 Wörter
\end{otherlanguage}

\printornamentstrue
\cleardoublepage

  
  % preface
  \addchap{Preface}

Before I start, I want to thank my advisor Prof.\ Dr.\ Dirk Pflüger.
It is his ideas and his valuable input that have driven me in
my time as a PhD student.
I am also grateful for the exciting time with the whole group of
SSE (Simulation Software Engineering) and
SGS (Simulation of Large Systems),
for which I want to thank all past and current PhD students and postdocs of
Dirk Pflüger and Prof.\ Dr.\ Miriam Mehl.
I thank
Carolin Schober,
Henriette Röger, and
Ralf Diestelkämper
for throughly proofreading
drafts of this thesis.

Similarly, I thank my co-advisor Prof.\ Oliver Röhrle, PhD, for the
interesting biomechanical collaboration and his willingness to be my co-advisor.
Many thanks go to the other members of the examination board for their
willingness to examine my thesis,
namely \todo{insert co-examiner} and \todo{insert co-examiner}.

Probably the most important role for the success of my PhD thesis
has played my family.
Without their mental support and distraction from the daily work,
I doubt that this thesis would have been possible.

Likewise, I am very grateful for the financial support from
the \foreignlanguage{ngerman}{Juniorprofessurenprogramm} of the
\foreignlanguage{ngerman}{Landesstiftung Baden-Württemberg}.
I thank the SimTech Cluster of Excellence for supporting
my three-month research stay in Canberra, Australia.

In addition, I want to thank the open-source community for making it possible to
write this thesis in an aesthetically sophisticated manner.
The list of software that was used to write this thesis includes
\LaTeX, Lua\LaTeX, Bib\LaTeX,
\scalebox{0.9}{\KOMAScript}, Ti\emph{k}Z, Python, Matplotlib,
and many more.

\label{page:preface}
Now, I wish that you, dear reader, obtain as much insight as possible
while reading the remaining
\pagedifference{page:preface}{LastPage} pages of this thesis.

Enjoy!

\vspace{1em}

\noindent
Stuttgart, \thedate

\noindent
\theauthor

  
  % TODO: dummy chapters
  \chapter{A Long Heading Without Any Meaning and Using Two Lines}
  Hello World!
  
  Umlaute: äöüß
    
  \begin{gather}
    X \times Y\\
    A \cdot \vec{x} = \vec{b}\\
    \min_{\vec{x} \in [0, 1]} \int_\Omega f(\vec{x}, \vec{y}) d\vec{y}
  \end{gather}
  
  \section{Abschnitt}
  
  \blindtext
  
  \subsection{Unterabschnitt}
  
  \blindtext
  
  \chapter{Kapitel 1}
  \blindmathpaper
  
  \chapter{Kapitel 2}
  \blindmathpaper
  
  % start appendix
  \appendix
  
  % TODO: cite all for debug purposes
  \nocite{*}
  % references (add to TOC)
  \printbibliography[heading=bibintoc]
\end{document}

\endinput









documentclass[
  paper=a4,
  twoside,
  fontsize=12pt,
  DIV=calc,
  headings=small,
  bibliography=totoc,
  listof=totoc,
  headsepline,
  footsepline,
  plainfootsepline,
]{scrbook}

%input{template}
usepackage{titlepage}

input{commands}
bibliography{../bib/bibliography}

\graphicspath{{../gfx/}}

% nomenclature
\renewcommand{\nomgroup}[1]{
  \ifthenelse{\equal{#1}{S}}{\item[\bfseries\sffamily\large{List of Symbols}]}{
    \ifthenelse{\equal{#1}{A}}{\item[\bfseries\sffamily\large{List of Abbreviations}]}{}}
}
\makenomenclature

\begin{document}
% hyphenations
input{hyphenation}

% Die Seitennummerierung erfolgt durchlaufend ab der Titelseite. Also keine
% Spielereien mit römischen Ziffern usw. - Die ISO 7144 schreibt das sogar für
% wissenschaftliche Werke vor.
% Von Promitionsordnung verlangt!
% Deshalb ist \frontmatter DEAKTIVIERT
%\frontmatter
\title{%
  B-Splines for Sparse Grids: Algorithms and
  Application to Higher-Dimensional Optimization%
}
\author{Julian Valentin}
\date{TODO}
\keywords{TODO}
\firstexaminer{Prof.~Dr.~Dirk Pflüger}
\secondexaminer{TODO\\TODO\\TODO}
\dateofexamination{TODO}
\placeofbirth{Stuttgart}
\faculty[Vom]{Stuttgart Research Centre for Simulation Technology}
\institute{Institut für Parallele und Verteilte Systeme}

\maketitle

% This is necessary if you have more than 9 sections/subsections/subsubsections
\makeatletter
\renewcommand\l@section{\@dottedtocline{1}{1.5em}{3em}}
\renewcommand\l@subsection{\@dottedtocline{2}{1.5em}{4.3em}}
\renewcommand\l@subsubsection{\@dottedtocline{3}{1.5em}{5.6em}}
\makeatother

\tableofcontents
\clearpage

\pagestyle{justpagenums}
% input{abstract}

%\mainmatter
\pagestyle{scrheadings}

%% ---- abstracts ----------------------------------------------
input{abstracts-danksagung}
%% -------------------------------------------------------------
\pagestyle{scrheadings}

%% ---- Content ------------------------------------------------
% input{examples}
input{introduction}
input{this_and_that}
%% -------------------------------------------------------------

\printbibliography
\noindent
Alle URLs wurden zuletzt am \today~gepr\"uft.

\clearpage
\listoffigures
\listoftables

\markleft{List of Theorems}{}
\markright{List of Theorems}{}
\chapter*{List of Theorems}
\addcontentsline{toc}{chapter}{List of Theorems}
\listtheorems{definition}

\appendix

% 'Anhang' ins Inhaltsverzeichnis
\phantomsection
\addcontentsline{toc}{part}{Appendix}

\chapter{Various}

%input{Z-Anhang}

%% --------------------------------------------------------------
% nomenclature
input{nomenclature}

\IfDefined{printindex}{\printindex}
\IfDefined{printnomenclature}{
  \markleft{Nomenclature}{}
  \markright{Nomenclature}{}
  \printnomenclature
}

\end{document}
